\firstsection{Introduction}

\maketitle

Tree-structured data are widely used to represent hierarchical information and model complex processes, such as organizational hierarchies, query plan trees in relational databases, reasoning trees in search and inference tasks, and abstract syntax trees (ASTs) in program analysis. In many real-world scenarios these trees are not static. They evolve over time or across steps as rules are applied, states are updated, search progresses, or human edits are introduced, forming sequences of tree transformations. Analyzing tree evolution is important for understanding how hierarchical structures change, identifying critical transition steps, and diagnosing abnormal or inefficient processes.

Compared with tree sequence analysis that treats trees primarily as a series of snapshots, Tree Evolution analysis focuses not only on what a tree looks like at a given step, but also with how one tree transforms into the next and why these changes occur. In many applications, the final tree only reflects the result of an evolving process, while obscuring key rewrites, backtracking behaviors, or local irregularities along the way. For example, in query optimization, an initial logical plan may go through a series of rule-based rewrites and structural adjustments before becoming an executable plan. Examining this process can help analysts understand optimizer behavior and locate the steps that lead to incorrect results or performance degradation. Similarly, in reasoning and search tasks, tree expansion, pruning, and rollback reveal the exploration strategy of the system; in program evolution analysis, local modifications on ASTs and their propagated effects can help expose refactoring behavior and error sources. These examples suggest that trees should often be understood as evolving structures rather than isolated snapshots.

Despite its importance, visual analytics of Tree Evolution remains challenging in at least three aspects. First, prior studies have shown that even comparing relatively simple graphs can be difficult for analysts~\cite{}, because topological comparison often disrupts cognitive continuity. This challenge becomes even greater in Tree Evolution analysis, where analysts must not only identify structural differences but also understand how one tree transforms into the next. Second, local topological changes are often scattered across large unchanged contexts, making it difficult to highlight critical differences without losing the structural context needed for interpretation. Third, critical changes in long evolution sequences are often sparse, making it hard to summarize major trends while preserving key transition steps. 

Existing studies have explored related problems from several perspectives, such as domain-specific evolution analysis, dynamic tree visualization, structural comparision, and temporal visual analytics. Some approaches are designed for specific application domains. For example, QOVIS~\cite{} supports the analysis of query
optimization process by leveraging domain knowledge. While such systems are effective in their target settings like the optimization rules. They rely on strong domain assumptions and are difficult to generalize to other tasks. Other approaches use animation~\cite{} or dynamic tree views~\cite{} to present transitions between adjacent steps, which can help reveal local continuity but become less effective for reviewing long sequences, locating critical steps, or comparing sparse changes across time. Temporal visual analytics methods can summarize trends and high-level stage-wise patterns, but they often ignore the topological structure and fine-grained change semantics of the trees. 
%Overall, existing approaches still fall short of jointly supporting topology-aware comparison between adjacent trees, context-preserving representation of critical differences, and long-sequence change summarization. 总觉得这里有点不到位，为什么需要我们的工作还是不够突出

%To address this gap, we propose \sysname{}, a topology-aware visual analytics framework for Tree Evolution. \sysname{} supports analysis at three levels. At the modeling level, it estimates corresponding changes between adjacent trees, representing evolution as a sequence of analyzable local change units rather than a set of disconnected snapshots. At the structural comparison level, it extracts compact difference-aware substructures that preserve essential context while reducing visual complexity, helping analysts focus on critical topological changes. At the sequence level, it organizes and simplifies local differences over time to reveal salient intervals, key transition steps, and overall evolutionary trends. By combining change modeling, topology-aware difference visualization, and sequence summarization, \sysname{} supports multi-granularity analysis from local changes to global evolution patterns.



%To address this gap, we propose \sysname{}, a topology-aware visual analytics framework for the evolution of tree structures. \sysname{} first computes topological differences between adjacent trees and identifies changed nodes, representing evolution as a sequence of local structural changes rather than disconnected snapshots. Based on these differences, it extracts and organizes difference-aware substructures that preserve essential context while reducing visual complexity, enabling effective comparison of topological differences without lose general context. \sysname{} further traces the lifecycle of nodes throughout the evolution process to support localized temporal analysis of changes. Building on these local analyses, it simplifies long sequences based on the distribution of differences to reveal salient intervals, key transition steps, and overall evolutionary trends. In this way, \sysname{} supports multi-granularity analysis from local changes to global evolution patterns.

To address this gap, we propose \sysname{}, a topology-aware visual analytics framework for the evolution of tree structures. \sysname{} first computes topological differences between adjacent trees and identifies changed nodes, representing the evolution as a sequence of local structural changes rather than a set of disconnected snapshots. It then extracts and organizes compact difference-aware substructures that preserve essential context while reducing visual complexity, enabling effective comparison of critical topological differences. Furthermore, \sysname{} traces the lifecycle of nodes throughout the evolution process to track the development of important information over time. Building on these local analyses, \sysname{} simplifies long sequences based on the distribution of differences to reveal salient intervals, key transition steps, and overall evolutionary trends. In this way, Building on these techniques, \sysname{} supports multi-granularity analysis from local changes to global evolution patterns.

The contributions of this work are as follows:

\squishlist
\item We formulate design requirements of \sysname, which addresses the core challenges of Tree Evolution as a visual analytics problem, including topology-aware comparison, context-preserving difference analysis, and long-sequence summarization.

%We formulate Tree Evolution as a visual analytics problem that jointly involves topology-aware comparison, context-preserving difference analysis, and long-sequence summarization, and we characterize common analytical needs shared across multiple application scenarios.

\item We propose a configurable differencing method for adjacent tree structures to estimate corresponding changes and evolution processes, and further develop a compact difference-aware substructure extraction strategy that balances fidelity and cognitive load.

\item We design and implement \sysname{}, a multi-granularity visual analytics framework that supports local topology comparison, cross-step evolution tracing, and long-sequence summarization of key processes.

\item We demonstrate the effectiveness and practical value of \sysname{} through multiple case studies and user interviews in tasks involving tree evolution understanding, key change localization, and long-sequence analysis.

\squishend

The remaining parts of the paper are organized as follows:
